\documentclass[a4paper,11pt]{article}
\usepackage{amsmath, amssymb, amsthm}

\newtheorem{lemma}{Lemma}

\title{Lokales Austauschlemma für List-Scheduling}
\author{}
\date{}

\begin{document}

\maketitle

\section*{Problemstellung}

Gegeben sei das Scheduling-Problem auf $m$ identischen parallelen Maschinen.
Für eine Permutation $\pi$ wird eine List-Schedule erzeugt: Jobs werden in der
Reihenfolge $\pi$ betrachtet und jeweils der aktuell am wenigsten belasteten
Maschine zugewiesen.

Für jeden Job $j$ mit Bearbeitungszeit $p_j > 0$ sei $S_j^\pi$ seine Startzeit.
Wir betrachten die Zielfunktion
\[
G(\pi) := \sum_{j=1}^n p_j S_j^\pi.
\]

\section*{Lokales Austauschlemma}

\begin{lemma}[Lokales Austauschlemma]
Seien $x$ und $y$ zwei Jobs, die in einer Permutation $\pi$ direkt hintereinander
auftreten. Sei $t$ die kleinste aktuelle Maschinenlast unmittelbar vor der
Zuweisung dieser beiden Jobs und $t+\delta$ die zweitkleinste Last mit
$\delta \ge 0$.

Dann gilt für die beiden möglichen lokalen Reihenfolgen $xy$ und $yx$:

\begin{itemize}
    \item Falls $p_x \le \delta$ und $p_y \le \delta$, dann
    \[
    G(xy) = G(yx).
    \]

    \item Falls $p_x > \delta$ und $p_y > \delta$, dann
    \[
    G(xy) = G(yx).
    \]

    \item Falls $p_x \le \delta < p_y$, dann
    \[
    G(xy) - G(yx) = p_x (p_y - \delta) > 0.
    \]

    \item Falls $p_y \le \delta < p_x$, dann
    \[
    G(yx) - G(xy) = p_y (p_x - \delta) > 0.
    \]
\end{itemize}
\end{lemma}

\section*{Beweis}

Die Zielfunktion kann über die Endlasten der Maschinen dargestellt werden als
\[
G(\pi) = \frac{1}{2} \left( \sum_{i=1}^m L_i^2 - \sum_{j=1}^n p_j^2 \right),
\]
wobei $L_i$ die Endlast von Maschine $i$ ist. Da der zweite Term konstant ist,
genügt es, $\sum_i L_i^2$ zu vergleichen.

Wir betrachten nur die lokale Wirkung der beiden Jobs $x$ und $y$. Sei vor deren
Zuweisung die kleinste Last $t$ und die zweitkleinste $t+\delta$.

\subsection*{Fall 1: $p_x \le \delta$ und $p_y \le \delta$}

In beiden Reihenfolgen werden beide Jobs auf dieselbe Maschine gelegt:
\[
L = t + p_x + p_y.
\]
Die resultierenden Lasten sind identisch, also auch $\sum_i L_i^2$.

\subsection*{Fall 2: $p_x > \delta$ und $p_y > \delta$}

Der erste Job geht auf die kleinste Maschine (Last $t$), danach ist deren Last
größer als $t+\delta$, sodass der zweite Job auf die zweitkleinste Maschine geht.
Die Endlasten sind in beiden Fällen:
\[
(t + p_x, \; t+\delta + p_y)
\quad \text{bzw.} \quad
(t + p_y, \; t+\delta + p_x).
\]
Die Summe der Quadrate ist gleich:
\[
{(t+p_x)}^2 + {(t+\delta+p_y)}^2 = {(t+p_y)}^2 + {(t+\delta+p_x)}^2.
\]

\subsection*{Fall 3: $p_x \le \delta < p_y$}

Für $xy$:
\[
\text{beide Jobs auf derselben Maschine} \Rightarrow L = t + p_x + p_y.
\]

Für $yx$:
\[
y \text{ auf } t, \quad x \text{ auf } t+\delta.
\]

Die Differenz der Quadratsummen ergibt:
\[
\Delta = {(t+p_x+p_y)}^2 + {(t+\delta)}^2 - \big[{(t+p_y)}^2 + {(t+\delta+p_x)}^2\big].
\]

Nach Ausmultiplizieren und Kürzen erhält man:
\[
\Delta = 2 p_x (p_y - \delta).
\]

Daraus folgt:
\[
G(xy) - G(yx) = \frac{1}{2} \Delta = p_x (p_y - \delta) > 0.
\]

\subsection*{Fall 4: $p_y \le \delta < p_x$}

Analog ergibt sich:
\[
G(yx) - G(xy) = p_y (p_x - \delta) > 0.
\]

\subsection*{Schluss}

Alle Fälle sind vollständig analysiert, und die Behauptung folgt.
\hfill $\square$

\end{document}
